\documentclass[11pt,letterpaper]{article}
\usepackage[margin=1in]{geometry}
\usepackage[utf8]{inputenc}
\usepackage[T1]{fontenc}
\usepackage{amsmath,amssymb,amsthm,amsfonts}
\usepackage{mathtools}
\usepackage{booktabs}
\usepackage{hyperref}
\usepackage{enumitem}
\newcommand{\bbR}{\mathbb{R}}
\newcommand{\bbH}{\mathbb{H}}
\newcommand{\bbS}{\mathbb{S}}
\newcommand{\cM}{\mathcal{M}}
\newcommand{\Lap}{\Delta}
\newcommand{\dvol}{\,d\mathrm{vol}}
\newcommand{\SO}{\mathrm{SO}}
\DeclareMathOperator{\Vol}{Vol}
\theoremstyle{plain}
\newtheorem{theorem}{Theorem}[section]
\newtheorem{proposition}[theorem]{Proposition}
\newtheorem{lemma}[theorem]{Lemma}
\newtheorem{corollary}[theorem]{Corollary}
\theoremstyle{remark}
\newtheorem{remark}[theorem]{Remark}

\title{The $L^2$ spectral geometry of the charge-one $\mathrm{SU}(2)$ instanton moduli space on $S^4$}
\author{Vico Bonfioli\\\small\texttt{vicobonfioli@gmail.com}}
\date{\today}

\begin{document}
\maketitle

\begin{abstract}
Let $\cM_1(S^4_r)$ be the moduli space of charge-one anti-self-dual $\mathrm{SU}(2)$ connections on the round $4$-sphere of radius $r$, with its natural $L^2$ (Weil--Petersson) metric $g_{L^2}$. By Groisser--Parker and Doi--Matsumoto--Matumoto, $(\cM_1(S^4_r),g_{L^2})$ is the interior of a compact smooth Riemannian $5$-manifold with boundary: an $\SO(5)$-invariant, cohomogeneity-one metric on the closed $5$-ball whose boundary is a totally geodesic round $S^4$ (the Uhlenbeck bubbling locus) at finite distance, of finite volume and strictly positive, non-constant sectional curvature. The same underlying $5$-ball also carries the information (Fisher--Rao) metric, which is hyperbolic $\bbH^5$; the two metrics have different spectral type. We determine the $L^2$ spectral geometry. The Laplace--Beltrami operator is symmetric but not essentially self-adjoint on the open moduli space; its self-adjoint realizations are the elliptic boundary conditions on the bounding $S^4$, and every such realization has purely discrete spectrum and obeys Weyl's law, with leading constant computed explicitly. Separating by $\SO(5)$ angular momentum reduces the problem to an explicit family of regular--singular Sturm--Liouville problems. A Liouville transformation, elementary here because $p\,w=g^4$, puts these in Schr\"odinger form with a potential increasing in both the radial variable and the angular momentum, and comparison against piecewise constant minorants of that potential yields certified two-sided enclosures. In the Groisser--Parker normalization the Friedrichs (Dirichlet) ground state is simple, $\SO(5)$-invariant, and lies in $[0.8203116,\,0.8203177]\,r^{-2}$; the Neumann spectral gap lies in $[0.9988337,\,0.9988386]\,r^{-2}$ and is the five-fold degenerate $k=1$ mode, in the vector representation of $\SO(5)$ (the instanton translation mode). We prove that this multiplet, and not any radial mode, is the gap. (The absolute values carry the metric normalization; the $\SO(5)$ multiplet structure, the level ordering, and discreteness do not.) The boundary $4$-volume is exactly $\tfrac23\pi^6 r^4$. We also record the exact $L^2$ interaction of two separated bubbles, an ingredient for the higher-charge boundary geometry.
\end{abstract}

\section{Introduction}\label{sec:intro}

The moduli space $\cM_k(S^4_r)$ of charge-$k$ ASD $\mathrm{SU}(2)$ connections on $S^4_r$ carries a natural $L^2$ metric $g_{L^2}$, induced by the $L^2$ inner product on connection $1$-forms. For $k=1$ the underlying manifold is the open $5$-ball, and the geometry of $g_{L^2}$ was determined explicitly by Groisser--Parker \cite{GP1987} and, equivalently, by Doi--Matsumoto--Matumoto \cite{DMM1987}; see also Franchetti--Schroers \cite{FS2016}, who used it for the classical adiabatic dynamics of instantons on $S^4$. We take that metric as given and determine the associated spectral theory of the $L^2$-Laplacian, together with its organization under the isometry group $\SO(5)$ and its interpretation as the collective-coordinate quantum mechanics of a single instanton.

The manifold underlying $\cM_1(S^4_r)$ is the open $5$-ball, as is hyperbolic space $\bbH^5$; the information (Fisher--Rao) metric on $\cM_1$ is the hyperbolic metric \cite{Info}. The $L^2$ metric is a different $\SO(5)$-invariant metric on the same ball: incomplete, of finite volume, and of strictly positive curvature, where $\bbH^5$ is complete and of constant negative curvature. The two Laplacians accordingly have different spectral type: that of $\bbH^5$ is purely continuous, $[4/r^2,\infty)$, while that of $g_{L^2}$ is purely discrete. The $L^2$ spectrum is also the set of energy levels of the moduli-space (geodesic) approximation to a quantized single instanton, organized into $\SO(5)$ multiplets of dimensions $1,5,14,30,55,\dots$

Spectral theory for a Weil--Petersson metric on a moduli space is not new. Ji, Mazzeo, M\"uller and Vasy \cite{JMMV} treat the Weil--Petersson Laplacian on the Riemann moduli space of genus $\gamma>1$, proving essential self-adjointness, discreteness of the spectrum, and a Weyl law. The present setting is the instanton analogue, and differs in one structural respect. There the metric is singular, with crossing cusp-edge singularities along a normal crossing divisor, and the Laplacian is essentially self-adjoint, so the self-adjoint realization is unique. Here $(\cM_1(S^4_r),g_{L^2})$ completes to a \emph{smooth} compact manifold with boundary; the Laplacian is correspondingly \emph{not} essentially self-adjoint, and its realizations form the family of elliptic boundary conditions on the bubbling $S^4$ described in Theorem~\ref{thm:discrete}. The analysis is easier here for the same reason.

\paragraph{The metric.} In $\SO(5)$-invariant cohomogeneity-one coordinates, with $x\in(0,1]$ a scale coordinate and the round $S^4$ the orbit of centers,
\begin{equation}\label{eq:gp-metric}
  g_{L^2} \;=\; r^2\big(\, f(x)\,dx^2 \;+\; g(x)\, g_{S^4}\,\big),
\end{equation}
where $g_{S^4}$ is the unit round metric and, in the normalization of \cite{GP1987,DMM1987},
\begin{align}
  f(x) &= \frac{4\pi^2}{(x^2-1)^4}\Big(x^4+10x^2+1 - \tfrac{12x^2(x^2+1)}{x^2-1}\ln x\Big),\label{eq:f}\\
  g(x) &= \frac{\pi^2}{2(x^2-1)^2}\Big(x^4-8x^2+1 + \tfrac{24x^4}{x^4-1}\ln x\Big).\label{eq:g}
\end{align}
The point $x=1$ is a smooth interior point, the unique $\SO(5)$-invariant instanton (the center of the ball), where the $S^4$ orbit collapses with $\sqrt{g}\sim(1-x)$ and unit cap derivative (no cone angle). The limit $x\to0$ is the bubbling boundary: a round $S^4$ of radius $\sqrt{g(0)}\,r = \sqrt{\pi^2/2}\;r$ at finite $L^2$-distance, totally geodesic (the radius has vanishing normal derivative there). The metric is geodesically incomplete, with
\begin{equation}\label{eq:geom-constants}
  \Vol(\cM_1,g_{L^2}) \in [943.29,\,943.50]\,r^5,\quad
  \Vol(\partial) = \tfrac{2}{3}\pi^6 r^4,\quad
  \operatorname{diam_{rad}} = 3.7437\,r,
\end{equation}
and strictly positive, non-constant sectional curvature (isotropic value $\approx 0.2533\,r^{-2}$ at the center, decreasing toward the boundary), in agreement with \cite[Cor.~C]{GP1987}. The volume enclosure is obtained from monotone Riemann sums for $\int\sqrt f\,g^2$, whose integrand is decreasing; it is not a recognizable closed form. The boundary $4$-volume is exactly $\tfrac23\pi^6r^4$ because $g(0)=\pi^2/2$ exactly.

\begin{remark}[On a hyperbolic misidentification]\label{rmk:not-hyperbolic}
The hyperbolic metric on the same $5$-ball is the information (Fisher--Rao) metric \cite{Info}, \emph{not} $g_{L^2}$. Being incomplete and positively curved, $g_{L^2}$ has none of the continuous $[4/r^2,\infty)$ spectrum of $\bbH^5$; its spectrum is discrete (Theorem~\ref{thm:discrete}).
\end{remark}

\section{The spectral type of the $L^2$-Laplacian}\label{sec:spectral}

Let $-\Lap$ be the Laplace--Beltrami operator of $g_{L^2}$, initially on $C_c^\infty$ of the open moduli space (functions supported away from the bubbling boundary).

\begin{theorem}[Discrete spectral type and Weyl law]\label{thm:discrete}
The operator $-\Lap$ on $C_c^\infty(\cM_1(S^4_r))$ is symmetric and bounded below but not essentially self-adjoint; its self-adjoint extensions correspond to elliptic boundary conditions on the totally geodesic bounding $S^4$. Every such extension has compact resolvent, hence purely discrete spectrum $0\le\lambda_0\le\lambda_1\le\cdots\to\infty$, with leading Weyl asymptotics
\[
  N(\lambda)=\#\{n:\lambda_n\le\lambda\}
  = \frac{\omega_5}{(2\pi)^5}\Vol(\overline{\cM}_1)\,\lambda^{5/2}\,(1+o(1)),
\]
where $\omega_d$ is the unit $d$-ball volume; in the normalization \eqref{eq:gp-metric} at $r=1$ the constant is $\tfrac{\omega_5}{(2\pi)^5}\Vol = 0.50710$. There is no continuous spectrum.
\end{theorem}

\begin{proof}[Proof (sketch)]
The metric completion $\overline{\cM}_1$ is a compact smooth manifold with boundary \cite[Cor.~C]{GP1987}; $f,g$ in \eqref{eq:f}--\eqref{eq:g} extend smoothly with finite nonzero limits at $x=0$ ($f(0)=4\pi^2$, $g(0)=\pi^2/2$), and $x=1$ is a smooth removable cap. Thus $-\Lap$ is uniformly elliptic with smooth coefficients on a compact manifold with boundary. Failure of essential self-adjointness on the interior $C_c^\infty$ and the classification of extensions by boundary conditions on $\partial\overline{\cM}_1=S^4$ are the standard boundary theory. Compactness of the resolvent follows from the compact Sobolev embedding $H^1(\overline{\cM}_1)\hookrightarrow L^2(\overline{\cM}_1)$, giving discrete spectrum. The Weyl term is standard.
\end{proof}

\begin{remark}[On the second Weyl term]\label{rmk:weyl2}
The refinement $\mp\tfrac14\tfrac{\omega_4}{(2\pi)^4}\Vol(\partial)\,\lambda^{2}$ ($-$ Dirichlet, $+$ Neumann), whose constant is $0.50734$ in the normalization \eqref{eq:gp-metric}, holds under Ivrii's condition that the periodic billiard trajectories carry measure zero. We do not verify that condition for this metric, and it is not automatic here: the boundary is a totally geodesic round $S^4$, so the trajectories gliding along it are precisely the geodesics of a round sphere and are therefore all closed, though they form a measure zero set in the cosphere bundle. The interior flow is not analyzed.
\end{remark}

\section{$\SO(5)$ reduction and the low spectrum}\label{sec:reduction}

Because $g_{L^2}$ is $\SO(5)$-invariant of cohomogeneity one, $L^2(\cM_1)$ decomposes over the $\SO(5)$-harmonics on the orbit $S^4$. Writing an eigenfunction as $u(x)\,Y_k(\omega)$ with $Y_k$ a degree-$k$ spherical harmonic ($\Lap_{S^4}Y_k=-k(k+3)Y_k$, multiplicity $m_k=\tfrac{(2k+3)(k+1)(k+2)}{6}=1,5,14,30,55,\dots$), the eigenvalue equation $-\Lap(uY_k)=\lambda u Y_k$ becomes the regular--singular Sturm--Liouville problem
\begin{equation}\label{eq:SL}
  -\big(p\,u'\big)' + k(k+3)\,\sqrt{f}\,g\,u \;=\; \lambda\,w\,u,
  \qquad p=\frac{g^2}{\sqrt f},\quad w=\sqrt{f}\,g^2,\quad x\in(0,1),
\end{equation}
with regularity at the center $x=1$ (where $u\sim(1-x)^k$) and the chosen elliptic condition at the boundary $x=0$. We solve \eqref{eq:SL} in three ways. Two are linear finite element discretizations that bracket the spectrum numerically: the conforming, consistent-mass scheme is Rayleigh--Ritz and bounds every eigenvalue from above, while the lumped-mass scheme converges from below, the two agreeing to width below $10^{-7}$ at $2.4\times10^4$ elements. The coefficients \eqref{eq:f}--\eqref{eq:g} are only $C^1$ at $x=0$, carrying an $x^2\ln x$ term that the Neumann modes feel and the Dirichlet modes, vanishing there, do not.

The third method is certified. Because $p\,w=g^4$, the Liouville transformation is elementary: with arc length $s=\int\sqrt f\,dx$ and $v=g\,u$, \eqref{eq:SL} becomes
\begin{equation}\label{eq:schrod}
  -v'' + Q_k\,v = \lambda\,v \quad\text{on }(0,L),\qquad
  Q_k = \frac{k(k+3)}{g} + \frac{g_{ss}}{g},\qquad L=3.7437,
\end{equation}
with $v$ Dirichlet at the center, where $g\to0$, and inheriting the chosen condition at the bubbling boundary, where $g_s=0$ because that boundary is totally geodesic. Near the center $g\sim s^2$ and $Q_k\sim 2/s^2$, the expected five-dimensional radial barrier. The function $Q_k$ is increasing in $s$ and increasing in $k$. Replacing it on each cell of a mesh by its value at the left endpoint therefore gives a piecewise constant $Q_k^-\le Q_k$ exactly, and the resulting comparison problem, solvable in closed form cell by cell, bounds the eigenvalues from below by min--max. The deficit is first order in the mesh width, so the enclosures tighten linearly under refinement; those quoted below use $2\times10^5$ cells and have width below $10^{-5}$.

\begin{theorem}[Ground state, gap, and multiplet structure]\label{thm:groundstate}
For the Friedrichs (Dirichlet) extension $-\Lap_D$ the ground state is simple, $\SO(5)$-invariant, realized by a positive eigenfunction, with (all absolute values in the normalization \eqref{eq:gp-metric})
\[
  \lambda_0(-\Lap_D)\;\in\;[0.8203116,\,0.8203177]\;r^{-2}\quad(k=0).
\]
For the Neumann extension $-\Lap_N$ the ground state is $0$ (the constant, $k=0$), and the spectral gap is \emph{not} an $\SO(5)$-invariant mode: it is the five-fold degenerate $k=1$ mode
\[
  \lambda_1(-\Lap_N)\;\in\;[0.9988337,\,0.9988386]\;r^{-2}\quad(k=1,\ \text{multiplicity }5,\ \text{vector representation of }\SO(5)).
\]
Both enclosures are certified. That the gap is this multiplet is likewise proved, not merely computed: $Q_k$ in \eqref{eq:schrod} increases with $k$, so the eigenvalues increase with $k$ at each radial position, and the certified bounds
\[
  \lambda(k=0,\ \text{first overtone})\ \ge\ 2.0351863\,r^{-2},
  \qquad \lambda(k=2,\ \text{ground})\ \ge\ 2.3973341\,r^{-2}
\]
separate the $k=1$ multiplet from everything above it. The low eigenvalues $\lambda\,r^2$, as bracketed by the two finite element schemes, are:
\begin{center}\small
\begin{tabular}{@{}c c l l@{}}
\toprule
$k$ & $m_k$ & Dirichlet (Friedrichs), lowest $\lambda r^2$ & Neumann, lowest $\lambda r^2$\\
\midrule
$0$ & $1$  & $0.820318,\ 3.611916,\ 7.830235$ & $0,\ 2.035539,\ 5.544165$\\
$1$ & $5$  & $2.180655,\ 5.667069,\ 10.584265$ & $0.998839,\ 3.742777,\ 7.948320$\\
$2$ & $14$ & $3.938201,\ 8.125416,\ 13.743074$ & $2.397606,\ 5.851182$\\
$3$ & $30$ & $6.096526,\ 10.986855$ & $4.199088,\ 8.361619$\\
$4$ & $55$ & $8.657372$ & $6.404425$\\
\bottomrule
\end{tabular}
\end{center}
We make no claim about closed forms for these eigenvalues. At the precision available here, about eight digits, inverse-symbolic searches return spurious matches against several bases, and are evidence in neither direction. What can be said concerns the equation rather than its eigenvalues: the coefficients \eqref{eq:f}--\eqref{eq:g} carry logarithms, so \eqref{eq:SL} is not hypergeometric.
\end{theorem}

\begin{proof}[Proof (sketch)]
Simplicity and positivity of the Dirichlet ground state, and its $\SO(5)$-invariance, are the Perron--Frobenius property of the Friedrichs Laplacian on a connected domain together with uniqueness of the positive ground state and invariance of the operator; a positive ground state is $\SO(5)$-fixed, hence $k=0$. The Neumann ground state is the constant. The remaining statements are the numerical solution of \eqref{eq:SL} for $k=0,\dots,4$; ordering the union of the per-$k$ spectra puts the five-dimensional $k=1$ vector multiplet at $0.9988386\,r^{-2}$, strictly below the first $k=0$ overtone. Convergence data are in the ancillary computational files.
\end{proof}

\begin{remark}[Collective-coordinate quantum mechanics]\label{rmk:qm}
In the moduli-space approximation \cite{FS2016}, slow motion on $(\cM_1,g_{L^2})$ is geodesic, and the associated quantum Hamiltonian is $-\tfrac12\Lap_{L^2}$; the eigenvalues of Theorem~\ref{thm:groundstate} are its energy levels. One qualification is essential. Here $\cM_1$ is the \emph{unframed} moduli space, of dimension five (center and scale). The collective-coordinate quantization standard in the soliton literature retains the $\mathrm{SU}(2)$ framing, and it is the framing directions that carry the spin and isospin quantum numbers; those are absent here. The spectrum below is therefore that of the unframed reduction, not the full collective-coordinate spectrum. Within that reduction, the gap being the $k=1$ vector multiplet says that the lowest excitation is a rigid translation of the center rather than a change of scale, the scale (breathing) excitation being the first $k=0$ overtone, and the $\SO(5)$ multiplet dimensions $1,5,14,30,\dots$ are the degeneracies of the levels.
\end{remark}

\section{The two-bubble $L^2$ interaction (higher-charge boundary geometry)}\label{sec:cross-term}

For $k\ge2$ the $L^2$ metric is that of a compact manifold with corners, the Uhlenbeck--Donaldson compactification adding finite-distance strata $\cM_{k'}(S^4_r)$, $k'<k$; Theorem~\ref{thm:discrete} generalizes verbatim (compact metric completion $\Rightarrow$ compact resolvent) to purely discrete spectrum with the $(8k-3)/2$-dimensional Weyl law. The content for $k\ge2$ is the near-boundary geometry, governed by the $L^2$ interaction of separated bubbles, the instanton analog of the Gibbons--Manton asymptotic metric for monopoles. We record the leading interaction in closed form (as $L^2$ geometry, not spectrum).

\begin{proposition}[Leading scale interaction at the maximally-bubbled corner]\label{prop:corner}
Near the corner of $\overline{\cM}_k(S^4_r)$ where $k$ bubbles have small scales $s_1,\dots,s_k$ at distinct centers with pairwise chordal separations $R_{ij}$, the scale--scale block of $g_{L^2}$ is, to leading order,
\[
  \big(g_{L^2}\big)_{s_i s_j} = 16\pi^2\,\delta_{ij} + (1-\delta_{ij})\,\frac{48\pi^2 s_i s_j}{R_{ij}^2} + O\!\big((s/R)^4\big),
\]
the diagonal the scale-independent BPST self-norm and the off-diagonal the two-bubble interaction of Lemma~\ref{lem:cross-term}.
\end{proposition}

\begin{lemma}[Exact two-bubble scale interaction]\label{lem:cross-term}
Let $A_1,A_2$ be BPST instantons on $\bbR^4$ with scales $s_1,s_2$, centers at separation $R$. Then
\[
  \Big\langle \frac{\partial A_1}{\partial s_1},\frac{\partial A_2}{\partial s_2}\Big\rangle_{L^2(\bbR^4)}
  = 48\pi^2 s_1 s_2 \int_0^1 t(1-t)\,\frac{2X(t)-t(1-t)R^2}{X(t)^2}\,dt,
  \quad X(t)=t(1-t)R^2+(1-t)s_1^2+t s_2^2,
\]
which evaluates in elementary closed form with leading asymptotic $48\pi^2 s_1s_2/R^2$ and exact next-order coefficient $-1$. (Machine-verified to $\sim2.5$ ulp against the $4$-dimensional integral.)
\end{lemma}

\begin{remark}
The closed form is correct as an $L^2$ integral; its earlier interpretation as a function of a hyperbolic distance referred to the information metric and is not used here.
\end{remark}

\begin{thebibliography}{9}
\bibitem{GP1987} D.\ Groisser and T.~H.\ Parker, \emph{The Riemannian geometry of the Yang--Mills moduli space}, Comm.\ Math.\ Phys.\ \textbf{112} (1987) 663--689.
\bibitem{DMM1987} H.\ Doi, Y.\ Matsumoto, T.\ Matumoto, \emph{An explicit formula of the metric on the moduli space of BPST-instantons over $S^4$}, in \emph{A F\^ete of Topology}, Academic Press, 1987.
\bibitem{FS2016} G.\ Franchetti and B.~J.\ Schroers, \emph{Adiabatic dynamics of instantons on $S^4$}, Comm.\ Math.\ Phys.\ \textbf{353} (2017) 185--228; arXiv:1508.06566.
\bibitem{JMMV} L.\ Ji, R.\ Mazzeo, W.\ M\"uller, A.\ Vasy, \emph{Spectral theory for the Weil--Petersson Laplacian on the Riemann moduli space}, arXiv:1206.4010.
\bibitem{Info} D.\ Groisser and M.~K.\ Murray, \emph{Instantons and the information metric}, Ann.\ Global Anal.\ Geom.\ \textbf{15} (1997) 519--537; arXiv:dg-ga/9611008. For charge one the information metric is the hyperbolic metric on the $5$-ball.
\end{thebibliography}

\end{document}
